% Preamble for the Moodle handout (format: plate-pdf). The job here is to make a
% LaTeX document look like it belongs to the deck: same typeface, same colours,
% same restraint about rules and weight. Nothing below is Reveal-aware — this
% file only ever sees the linear document Pandoc builds from the same source.
%
% Two things arrive from elsewhere and are deliberately not hardcoded here:
%
%   * the brand colours, injected as \definecolor by handout.lua so _brand.yml
%     stays the single source of truth;
%   * the fonts, loaded by relative path from the extension.
%
% See the notes at each site below.

% ------------------------------------------------------------------ colours
%
% Defaults, in case handout.lua did not run (rendering this format outside the
% project, say). Guarded so that whichever of the two blocks LaTeX sees first,
% the brand wins: if the filter already defined a colour, this leaves it alone;
% if the filter's block comes later, its \definecolor simply overrides.
\makeatletter
\@ifundefined{\string\color@plate-ink}{\definecolor{plate-ink}{HTML}{15181B}}{}
\@ifundefined{\string\color@plate-primary}{\definecolor{plate-primary}{HTML}{2A5B8C}}{}
\@ifundefined{\string\color@plate-danger}{\definecolor{plate-danger}{HTML}{A6323F}}{}
\@ifundefined{\string\color@plate-muted}{\definecolor{plate-muted}{HTML}{6E747C}}{}
\makeatother

% The hairline used for the header rule and under tables. Mixed rather than
% pinned, so it tracks the ink colour like $rule does in the SCSS.
\colorlet{plate-rule}{plate-ink!22!white}

% -------------------------------------------------------------------- fonts
%
% The same Inter the deck uses, unpacked from woff2 to ttf by
% fonts/build-pdf-fonts.sh. The path is relative because every deck in this
% project sits exactly one folder below the root (week-01/index.qmd,
% template/index.qmd) and LaTeX runs in the deck's own folder. A deck nested any
% deeper would need this adjusted — that is the one assumption this file makes
% about where it is called from.
%
% Light (300) for headings is the whole look of the theme; Regular for body,
% because Light at 10.5pt on paper is too pale to read comfortably. That split is
% the one deliberate difference from the slides, where everything is Light.
\usepackage{fontspec}
\defaultfontfeatures{Ligatures=TeX}

\newfontfamily\plateLight{Inter-Light.ttf}[
  Path = ../_extensions/adamcoates/plate/fonts/ttf/,
  ItalicFont = Inter-LightItalic.ttf,
]

\setmainfont{Inter-Regular.ttf}[
  Path = ../_extensions/adamcoates/plate/fonts/ttf/,
  ItalicFont = Inter-Italic.ttf,
  BoldFont = Inter-SemiBold.ttf,
  BoldItalicFont = Inter-Italic.ttf,
  Numbers = {Lining, Proportional},
]

\setsansfont{Inter-Regular.ttf}[
  Path = ../_extensions/adamcoates/plate/fonts/ttf/,
  ItalicFont = Inter-Italic.ttf,
  BoldFont = Inter-SemiBold.ttf,
]

% Emoji. The decks use them as inline markers (📧 before an address, 🏛️ before an
% office). Inter has no such glyphs, so LaTeX drops them silently — the browser
% renders them fine, so the gap only ever shows up in the PDF.
%
% handout.lua rewrites the ones that appear in these decks into FontAwesome
% glyphs, which need the package loaded whether or not the document also has a
% callout (Quarto pulls fontawesome5 in only when one is present).
%
% A colour-emoji fallback font was the other option and is not worth it: Noto
% Color Emoji is a bitmap face that lualatex renders poorly, and it is invisible
% inside the nix devShell in any case.
\usepackage{fontawesome5}

% ----------------------------------------------------------------- headings
%
% Mirrors the deck: large and Light with space under it, never a rule or a
% coloured bar. h3 is the accented sub-heading (italic lapis) and h4 the warning
% register (claret) — the same two flourishes the slides allow themselves.
\addtokomafont{disposition}{\plateLight}
\setkomafont{section}{\plateLight\LARGE\color{plate-ink}}
\setkomafont{subsection}{\plateLight\Large\color{plate-ink}}
\setkomafont{subsubsection}{\plateLight\large\itshape\color{plate-primary}}
\setkomafont{paragraph}{\normalsize\color{plate-danger}}

% One page per slide. In the deck `#` opens a section and `##` is a slide; here
% they become \section and \subsection, so a break before each keeps the handout
% and the projector on the same page — literally. \clearpage rather than
% \newpage so a figure cannot float past the slide it belongs to and land under
% the next slide's heading.
%
% etoolbox's \pretocmd rather than titlesec's \sectionbreak: titlesec and KOMA
% each want to own sectioning and warn loudly when both are loaded.
%
% Deferred to \AtBeginDocument, and that is not caution — it is required.
% \RedeclareSectionCommand at the foot of this file *rebuilds* \section, which
% throws away any patch applied to it beforehand. Patching once the preamble has
% finished is the only point at which nothing can redefine the command back.
% (\pretocmd reports success either way, so getting this wrong fails silently:
% the handout simply runs its slides together mid-page.)
\usepackage{tikz}
\usetikzlibrary{arrows.meta}
\usepackage{etoolbox}
\AtBeginDocument{%
  \pretocmd{\section}{\clearpage}{}{\PackageWarning{plate}{could not hook \string\section}}%
  \pretocmd{\subsection}{\clearpage}{}{\PackageWarning{plate}{could not hook \string\subsection}}%
}

% ------------------------------------------------------- header and footer
%
% A hairline, the title on the left, the page number on the right, and the
% course logo small on the right of the header. Quiet enough to ignore while
% reading, present enough that a printed page is identifiable.
\usepackage[headsepline]{scrlayer-scrpage}
\clearpairofpagestyles
\setkomafont{pageheadfoot}{\footnotesize\color{plate-muted}}
\setkomafont{pagenumber}{\footnotesize\color{plate-muted}}
\renewcommand*{\headfont}{\normalfont\footnotesize\color{plate-muted}}

% The rule under the header, in the same hairline grey the tables use.
\addtokomafont{headsepline}{\color{plate-rule}}
\setlength{\headheight}{16pt}

\ihead{\platehandoutshorttitle}
\ohead{\platehandoutlogo}
\cfoot{\pagemark}

% Both are filled in by handout.lua (the logo comes from _brand.yml, the short
% title from the deck's own YAML). Defaults keep a render without the filter from
% dying rather than from looking right.
\providecommand{\platehandoutshorttitle}{}
\providecommand{\platehandoutlogo}{}

% The first page of a KOMA article uses `plain`, which would otherwise drop the
% header entirely; keep it consistent with every other page.
\renewcommand*{\titlepagestyle}{scrheadings}

% ------------------------------------------------------------------- tables
%
% booktabs is what Pandoc emits for pipe tables. The deck's tables carry rules
% only under the header and under the last row, at hairline weight — match that
% rather than LaTeX's heavier default.
\usepackage{booktabs}
% colortbl is what defines \arrayrulecolor. Quarto loads xcolor but without the
% `table` option, so the rule-colouring commands are not there by default.
\usepackage{colortbl}
\setlength{\heavyrulewidth}{0.6pt}
\setlength{\lightrulewidth}{0.4pt}
\arrayrulecolor{plate-rule}
\renewcommand{\arraystretch}{1.15}

% ----------------------------------------------------------------- captions
%
% Quarto centres captions above tables and under figures. The deck ranges them
% left in small caps and quiet grey (see `.reveal table caption` in plate.scss);
% match that, or a centred line over a left-aligned table reads as a mistake.
\usepackage{caption}
\captionsetup{
  justification = raggedright,
  singlelinecheck = false,
  font = {small, color = plate-muted},
  labelfont = {small, color = plate-muted},
  skip = 4pt,
}

% ------------------------------------------------------------- justification
%
% Ranged left, like the deck. On a 16 cm measure full justification opens rivers
% around the long uni-graz URLs these decks are full of — LaTeX cannot hyphenate
% a URL, so it stretches the whole line instead. \RaggedRight keeps hyphenation
% on for ordinary words, unlike plain \raggedright.
\usepackage{ragged2e}
\AtBeginDocument{\RaggedRight}
\setlength{\RaggedRightRightskip}{0pt plus 2em}

% -------------------------------------------------------------------- links
%
% colorlinks is set in _extension.yml; this is only the colour. Everything
% clickable is lapis, matching the deck, and nothing is boxed.
%
% Deferred to \AtBeginDocument because Quarto's template loads hyperref *after*
% include-in-header — calling \hypersetup here directly is an undefined control
% sequence.
\AtBeginDocument{%
  \hypersetup{
    linkcolor = plate-primary,
    urlcolor  = plate-primary,
    citecolor = plate-primary,
    filecolor = plate-primary,
  }%
}

% --------------------------------------------------------------- title block
%
% Echoes the cover slide: title large and Light, subtitle italic lapis under it.
\setkomafont{title}{\plateLight\Huge\color{plate-ink}}
\setkomafont{subtitle}{\plateLight\Large\itshape\color{plate-primary}}
\setkomafont{author}{\normalsize\color{plate-ink}}
\setkomafont{date}{\normalsize\color{plate-muted}}

% ---------------------------------------------------------------------- toc
%
% One line per slide is a useful index for a student scanning the handout, but it
% wants to read as a list, not as a numbered outline.
\setkomafont{sectionentry}{\plateLight\large\color{plate-ink}}
\RedeclareSectionCommand[tocindent=0pt]{section}
